\clearpage
\section{Wnioski}
\label{sec:wnioski}
\begin{enumerate}
 \item Dla pięciu próbek przy stacji gazowej uzyskano średni
 energetyczny poziom $52{,}74\,\mathrm{dB}$, a dla pięciu próbek
 tła $48{,}08\,\mathrm{dB}$. Różnica wyniosła $4{,}66\,\mathrm{dB}$,
 dlatego wpływ tła wymagał uwzględnienia.
 \item Po korekcji tła otrzymano robocze oszacowanie poziomu
 związanego z pracą stacji $L_{Aek}\approx50{,}9\,\mathrm{dB}$.
 Wynik charakteryzuje wybrany punkt i warunki krótkiej serii
 pomiarowej.
 \item Seria przy stacji spełnia kryterium liczby próbek z instrukcji.
 Dla serii tła wymagane było siedem próbek, a dostępnych jest pięć.
 Jest to niezrealizowany element procedury.
 \item Nie przeprowadzono sprawdzenia toru kalibratorem przed
 sesją ani po niej, ponieważ nie był on dostępny. Brak dokumentacji
 wzorcowania i danych meteorologicznych ogranicza ocenę
 wiarygodności wyników.
\end{enumerate}
